\documentclass[12pt,a4paper]{article}

\usepackage[a4paper, top=2.5cm, bottom=2.5cm, left=2.8cm, right=2.8cm]{geometry}
\usepackage[T1]{fontenc}
\usepackage[utf8]{inputenc}
\usepackage{parskip}
\usepackage{setspace}
\usepackage{hyperref}
\usepackage{listings}
\usepackage{amsmath}
\usepackage{booktabs}
\usepackage{longtable}
\usepackage{enumitem}
\usepackage{graphicx}
\usepackage{caption}

\hypersetup{
    colorlinks=false,
    linkcolor=black,
    urlcolor=black,
}

\setstretch{1.3}
\setlength{\parindent}{0pt}
\setlength{\parskip}{8pt}

\lstset{
  basicstyle=\ttfamily\small,
  breaklines=true,
  frame=single,
  tabsize=4,
  showstringspaces=false,
  xleftmargin=0.5em,
  xrightmargin=0.5em,
}

% -----------------------------------------------------------
\begin{document}

% Title block
\begin{center}
  {\Large \textbf{AI-Based Village Pond Planning System}} \\[0.5em]
  {\large Assignment 1 Report} \\[0.3em]
  Computer System Design \\[0.3em]
  Avinash \quad | \quad September 2026 \\[0.3em]
  \textbf{GitHub:} \url{https://github.com/avinash979309/village-pond-planning-system}
\end{center}

\vspace{0.5em}
\hrule
\vspace{1em}

% -----------------------------------------------------------
\section*{1. Introduction}

The goal of this assignment was to build a backend API that can automatically
identify the best location for a village pond given a contour map of an area.
Village ponds are important for rural water storage, and picking the right site
matters a lot -- it needs to be at a low-lying spot that collects rainwater
from a large area, and should not be placed on top of existing rivers or canals.

Instead of doing this manually, the system takes a KML contour map as input,
runs terrain and hydrology analysis, and gives back the best pond location
along with the catchment area (the region that drains into that point).

The API is built in Python using FastAPI and can be tested with a simple
curl command or through the browser.

% -----------------------------------------------------------
\section*{2. GitHub Repository}

\textbf{Link:} \url{https://github.com/avinash979309/village-pond-planning-system}

The repository has the full source code, the sample contour map, and a helper
shell script to run the analysis from the command line.

\textbf{Folder structure:}

\begin{lstlisting}
village-pond-planning-system/
|-- backend/
|   |-- app/
|   |   |-- main.py                          # entry point + /analyzeContour
|   |   |-- services/
|   |   |   `-- contour_analysis_service.py  # orchestrates the pipeline
|   |   `-- geo/
|   |       |-- kml_parser.py
|   |       |-- terrain_builder.py
|   |       |-- hydrology_engine.py
|   |       |-- osm_water_fetcher.py
|   |       |-- pond_candidate_selector.py
|   |       `-- catchment_delineator.py
|   `-- requirements.txt
|-- maps/
|   `-- sample_contour_map.kml    # sample input (Shivnath river area)
|-- generate_results.sh           # saves output as result.geojson
|-- result.geojson
`-- README.md
\end{lstlisting}

% -----------------------------------------------------------
\section*{3. Working API Route}

\textbf{How to run the server:}

\begin{lstlisting}[language=bash]
cd backend
python3 -m venv venv
source venv/bin/activate
pip install -r requirements.txt
uvicorn app.main:app --port 8000
\end{lstlisting}

The server will start at \texttt{http://localhost:8000}. There are three routes:

\begin{center}
\begin{tabular}{lll}
\toprule
Method & Endpoint & What it does \\
\midrule
GET  & \texttt{/}              & Returns API info \\
GET  & \texttt{/api/v1/health} & Health check \\
POST & \texttt{/analyzeContour}& Main analysis (upload KML, get result) \\
\bottomrule
\end{tabular}
\end{center}

\textbf{Testing the main route:}

\begin{lstlisting}[language=bash]
# curl (single line, no backslash)
curl -X POST http://localhost:8000/analyzeContour \
  -F "file=@maps/sample_contour_map.kml"

# OR use the helper script
./generate_results.sh maps/sample_contour_map.kml
\end{lstlisting}

You can also use the Swagger UI at \texttt{http://localhost:8000/docs} -- it
lets you upload a file directly from the browser without any curl commands.

\textbf{Example response:}

\begin{lstlisting}
{
  "status": "success",
  "pond_location": {
    "longitude": 81.2944,
    "latitude": 21.2439,
    "elevation_m": 282.0,
    "suitability_score": 0.8822
  },
  "pour_point": {
    "longitude": 81.2951,
    "latitude": 21.2441
  },
  "catchment": {
    "area_km2": 0.0139,
    "area_m2": 13851.0,
    "avg_elevation_m": 290.4,
    "boundary_geojson": { "type": "Polygon", "coordinates": ["..."] }
  },
  "all_candidates": [
    { "rank": 1, "suitability_score": 0.8822, "catchment_area_km2": 0.0139 },
    { "rank": 2, "suitability_score": 0.8800, "catchment_area_km2": 0.0004 }
  ],
  "osm_water_exclusion": {
    "water_bodies_found": true,
    "water_body_count": 11,
    "water_body_names": ["Shivnath", "sivnath river", "canal"]
  }
}
\end{lstlisting}

% -----------------------------------------------------------
\section*{4. Catchment Estimation Approach}

The system follows a standard hydrology pipeline to estimate the catchment
area for each candidate pond site. The steps are explained below.

\subsection*{4.1 Parsing the KML file}

The uploaded KML file is read using the \texttt{lxml} library. Each contour
line is a \texttt{<Placemark>} inside the KML. The parser pulls out the
vertex coordinates (longitude, latitude) and the elevation of each line from
the placemark name. It also calculates the bounding box of the whole area.

\subsection*{4.2 Building the Elevation Grid (DEM)}

Contour lines only give elevation at certain positions. To do hydrology
analysis we need elevation at every point on a grid. So the system takes all
the contour vertices as scattered $(x, y, z)$ points and uses
\textbf{SciPy linear interpolation} (\texttt{scipy.interpolate.griddata})
to fill in a regular $200 \times 200$ grid. Any leftover empty cells are
filled using nearest-neighbour interpolation.

The result is a Digital Elevation Model (DEM) -- a 2D array where each cell
has an elevation value in metres.

\subsection*{4.3 Conditioning the DEM}

Raw DEMs often have ``sinks'' -- cells that are lower than all their
neighbours with no way for water to flow out. These break the flow analysis.
The library \textbf{pysheds} is used to:

\begin{itemize}[leftmargin=2em]
  \item \textbf{Fill sinks} -- raise the elevation of trapped cells until
        water can flow out.
  \item \textbf{Resolve flats} -- add a tiny slope to completely flat areas
        so the flow direction is not ambiguous.
\end{itemize}

\subsection*{4.4 D8 Flow Direction and Accumulation}

The \textbf{D8 algorithm} is used to model how water flows across the terrain.
For each cell, it looks at all 8 neighbours and sends the flow to the one
that is steepest downhill:

\[
\text{flow direction of cell}_{ij} = \arg\max_{k \in \{1 \ldots 8\}}
\left( \frac{z_{ij} - z_k}{L_k} \right)
\]

Here $z_{ij}$ is the elevation of the cell, $z_k$ is the neighbour's
elevation, and $L_k$ is the distance (1 for up/down/left/right, $\sqrt{2}$
for diagonal).

After flow directions are assigned, \textbf{flow accumulation} is calculated --
each cell gets a count of how many upstream cells drain through it. Cells
with very high accumulation are major drainage channels (rivers, streams).
The top 2\% are flagged and excluded from pond site selection.

\subsection*{4.5 Excluding Existing Water Bodies}

Before picking pond sites, the system removes locations that already have
water. Two methods are used:

\textbf{OpenStreetMap (OSM):} The Overpass API is queried to download polygons
of all rivers, lakes, and canals in the study area. Candidate cells inside
these polygons are excluded. The query uses multiple mirror servers in case
one is down.

\textbf{Terrain-based fallback:} If OSM is unreachable, cells that are both
very low-lying and have high flow accumulation are removed as likely river
corridors.

\subsection*{4.6 Scoring and Selecting Candidates}

Each remaining cell gets a score from 0 to 1 based on four things:

\begin{center}
\begin{tabular}{llr}
\toprule
Factor & Reason & Weight \\
\midrule
Low elevation relative to surroundings & Water naturally collects here & 30\% \\
Gentle slope                           & Stable site for construction  & 25\% \\
High flow accumulation                 & Large area drains into it     & 35\% \\
Distance from rivers/exclusion zones   & Safety buffer                 & 10\% \\
\bottomrule
\end{tabular}
\end{center}

The top 3 candidates are picked, with a minimum spatial separation between
them so they are not all in the same spot.

\subsection*{4.7 Delineating the Catchment}

Once the best site is chosen, the catchment is computed:

\begin{enumerate}[leftmargin=2em]
  \item \textbf{Snap to pour point} -- the candidate is moved to the nearest
        high-accumulation cell within a small search radius. This ensures it
        sits exactly on the drainage channel.
  \item \textbf{Trace upstream} -- pysheds follows the D8 flow directions
        backwards from the pour point, marking every cell that eventually
        drains into it. This set of cells is the catchment.
  \item \textbf{Calculate area} -- the number of catchment cells is multiplied
        by the area of one cell:
        \[
        A_{\text{cell}} =
          \frac{\Delta\text{lon} \cdot 111{,}320 \cdot \cos(\phi)}{N}
          \times
          \frac{\Delta\text{lat} \cdot 110{,}574}{N}
          \quad [\text{m}^2]
        \]
        where $\phi$ is the mean latitude and $N = 200$ is the grid size.
  \item \textbf{Get boundary} -- the raster mask is converted to a GeoJSON
        Polygon using Shapely, which is returned in the API response.
\end{enumerate}

% -----------------------------------------------------------
\section*{5. Demonstration using the Sample Contour Map}

The sample map provided (\texttt{maps/sample\_contour\_map.kml}) covers a
roughly 1 km\textsuperscript{2} area in the Shivnath river region in central
India (around $21.24°\,\text{N},\ 81.29°\,\text{E}$). It contains multiple
closed contour lines and the Shivnath river and some canals pass through the
lower part of the area.

\textbf{Command used:}
\begin{lstlisting}[language=bash]
./generate_results.sh maps/sample_contour_map.kml
\end{lstlisting}

\textbf{Result:}

\begin{center}
\begin{tabular}{ll}
\toprule
Field & Value \\
\midrule
Pond Location (Rank 1) & $81.2944°\,\text{E},\ 21.2439°\,\text{N}$ \\
Elevation              & 282.0 m \\
Suitability Score      & 0.8822 / 1.0 \\
Pour Point             & $81.2951°\,\text{E},\ 21.2441°\,\text{N}$ \\
Catchment Area         & 0.0139 km\textsuperscript{2} (13,851 m\textsuperscript{2}) \\
Avg Catchment Elevation& 290.4 m \\
Candidates found       & 2 \\
OSM features excluded  & Shivnath, sivnath river, canal (11 total) \\
\bottomrule
\end{tabular}
\end{center}

The recommended location is slightly lower than its surroundings at 282 m,
which makes sense as a water collection point. The catchment of
$\approx$0.014 km\textsuperscript{2} means about 14,000 m\textsuperscript{2}
of land drains into this spot -- around 1.4\% of the whole map area, which is
a reasonable contributing area for a small village pond.

The OSM layer correctly excluded the Shivnath river and the canals, so the
recommended site is not on or near an existing waterway. The suitability
score of 0.88 shows it is a good site overall.

To visualize: drag the generated \texttt{result.geojson} file into
\href{https://geojson.io}{\texttt{geojson.io}} to see the pond location and
catchment polygon on the map.

% -----------------------------------------------------------
\section*{6. API Documentation}

\subsection*{Parameters for \texttt{POST /analyzeContour}}

\begin{center}
\begin{tabular}{llll}
\toprule
Parameter & Type & Default & Description \\
\midrule
\texttt{file}                    & file  & --    & KML or KMZ file (required) \\
\texttt{grid\_resolution}        & int   & 200   & DEM grid size per axis \\
\texttt{drainage\_threshold\_pct}& float & 2.0   & Top-N\% flow cells = drainage \\
\texttt{drainage\_buffer\_cells} & int   & 2     & Buffer around drainage cells \\
\texttt{snap\_radius\_cells}     & int   & 5     & Pour-point snap search radius \\
\texttt{skip\_osm}               & bool  & false & Skip OSM (offline/faster mode) \\
\bottomrule
\end{tabular}
\end{center}

\subsection*{Error responses}

\begin{center}
\begin{tabular}{lll}
\toprule
Code & When & Response \\
\midrule
422 & File field missing or wrong type & \texttt{\{"detail": [...]\}} \\
400 & KML/KMZ cannot be parsed         & \texttt{\{"status": "error", ...\}} \\
500 & Analysis failed internally        & \texttt{\{"status": "error", ...\}} \\
\bottomrule
\end{tabular}
\end{center}

\subsection*{Tech stack}

\begin{center}
\begin{tabular}{ll}
\toprule
Library & Purpose \\
\midrule
FastAPI + Uvicorn & REST API and web server \\
lxml + Shapely    & KML parsing and geometry \\
NumPy + SciPy     & DEM interpolation \\
pysheds           & D8 flow, sink fill, catchment delineation \\
httpx             & Overpass API requests \\
\bottomrule
\end{tabular}
\end{center}

% -----------------------------------------------------------
\section*{7. Conclusion}

The system works as intended. Given a KML contour map, it automatically
identifies a suitable pond location by analysing terrain and drainage patterns,
excluding existing rivers using OpenStreetMap data, and returning a catchment
polygon and suitability score.

The demo on the Shivnath river area produces a geographically reasonable
result -- the chosen site avoids the river, sits at a natural low point, and
has a sensible catchment area.

All code is available at:\\
\url{https://github.com/avinash979309/village-pond-planning-system}

\end{document}
